\documentclass[10pt,a4paper]{amsart}
\usepackage[margin=19mm]{geometry}
\usepackage{amsmath,amssymb,mathtools}
\usepackage[scr=rsfs,cal=euler]{mathalfa}
\usepackage{xfrac}
\usepackage[unicode,hidelinks,bookmarks=false]{hyperref}
\newcommand{\ie}{i.e.\ }
\newcommand{\cf}{cf.\ }
\newcommand{\resp}{respectively\ }
\newcommand{\Subsection}{\subsection}
\providecommand{\nobreakdash}{\nobreak\hbox{-}}
\setlength{\emergencystretch}{2em}
\multlinegap0pt
\usepackage{amssymb}
\newtheorem{lemma}{Lemma}
\title{Normal Behavior and Periodic Points of a Pseudo-Aliquot Map Associated with the Dedekind Psi Function}
\author{Aimin Guo}
\date{Source: arXiv:2609.06332v1 (CC BY 4.0)}
\begin{document}
\maketitle

\noindent\emph{Source and licence.} Normal Behavior and Periodic Points of a Pseudo-Aliquot Map Associated with the Dedekind Psi Function. Version arXiv:2609.06332v1 (CC BY 4.0), distributed by arXiv under the Creative Commons Attribution 4.0 International licence, http://creativecommons.org/licenses/by/4.0/, as recorded in the arXiv licence metadata for this exact version and shown on the abstract page https://arxiv.org/abs/2609.06332v1. Full licence notice: \texttt{licenses/arxiv-2609.06332-CC-BY-4.0.txt}.\par\smallskip

\noindent\textit{Editorial scope.} Counted unit: the article's lemma lem:congruence-persistence (source0.tex lines 1298--1311). The article's complete original proof of it is delivered with it (source0.tex lines 1313--1476, one \texttt{proof} environment of 164 lines). No mathematics is written, repaired or re-derived: the statement and the proof are the article's own lines, in the article's own notation, and no retained line is substituted or re-flowed.  Retained uncounted as the source-local context the proof needs: lines 1264--1280.  Nothing was omitted from any retained span, so the package carries no omission record.  No retained line contains a citation, so the wrapper carries no bibliography and no bibliography entry.  No retained line contains a figure environment, an image-inclusion command or drawing code, so no image is delivered and the package declares no asset dependency.  Editorial formatting only, in addition: an AMS \texttt{amsart} preamble that restates the article's own declarations and macros used by the retained lines, namely the environments \texttt{lemma} on separate counters, together with the standard packages those lines need.  The retained environments print in this document as Lemma 1 in order of appearance, whereas the article prints the same items with the numbering of its own full document.  The complete native source of the article as distributed in the arXiv e-print for this exact version is bundled unchanged as \texttt{sources/209524/source0.tex}.
\begin{lemma}\label{lem:prime-chain}
Let $M$ and $\ell$ be fixed positive integers. Then, for almost all
positive integers $n$, there exist prime divisors
\[
q_1,\ldots,q_\ell
\]
of $n$ such that
\[
q_1\equiv-1\pmod M
\]
and
\[
q_{r+1}\equiv-1\pmod{q_r},
\qquad
1\le r<\ell.
\]
\end{lemma}

\begin{lemma}\label{lem:congruence-persistence}
Let $K$ and $M$ be fixed positive integers. Then, for almost all
positive integers $n$,
\[
T^{(j)}(n)>1
\]
and
\[
T^{(j)}(n)\equiv(-1)^jn\pmod M,
\qquad
0\le j\le K,
\]
simultaneously.
\end{lemma}

\begin{proof}
Apply Lemma~\ref{lem:prime-chain} with
\[
\ell=K+1.
\]
Outside a set of asymptotic density $0$, the integer $n$ has prime
divisors
\[
q_1,\ldots,q_{K+1}
\]
satisfying
\[
q_1\equiv-1\pmod M
\]
and
\[
q_{r+1}\equiv-1\pmod{q_r},
\qquad
1\le r\le K.
\]

We first show simultaneously that
\[
T^{(j)}(n)>1
\]
and
\[
q_1,\ldots,q_{K+1-j}\mid T^{(j)}(n),
\qquad
0\le j\le K.
\]
For $j=0$, we have
\[
T^{(0)}(n)=n.
\]
Since every $q_r$ divides $n$, and in particular $q_1\mid n$, we
have
\[
n\ge q_1\ge2.
\]
Thus the assertion holds for $j=0$.

Suppose that the assertion holds for some $j<K$. Then
\[
T^{(j)}(n)>1.
\]
Let
\[
1\le r\le K-j.
\]
By the induction hypothesis,
\[
q_r\mid T^{(j)}(n)
\qquad\text{and}\qquad
q_{r+1}\mid T^{(j)}(n).
\]
Write
\[
q_{r+1}^a\parallel T^{(j)}(n)
\]
for some $a\ge1$. Since
\[
\psi(q_{r+1}^a)
=
q_{r+1}^{a-1}(q_{r+1}+1)
\]
and
\[
q_r\mid q_{r+1}+1,
\]
the multiplicativity of $\psi$ gives
\[
q_r\mid\psi\bigl(T^{(j)}(n)\bigr).
\]
Since also
\[
q_r\mid T^{(j)}(n),
\]
we obtain
\[
q_r
\mid
\psi\bigl(T^{(j)}(n)\bigr)-T^{(j)}(n)
=
T^{(j+1)}(n).
\]
Hence
\[
q_1,\ldots,q_{K-j}\mid T^{(j+1)}(n).
\]

Moreover, $T(m)>0$ for every integer $m>1$, since
\[
T(m)
=
m\left(
\prod_{p\mid m}\left(1+\frac1p\right)-1
\right)>0.
\]
Thus
\[
T^{(j+1)}(n)>0.
\]
Since
\[
q_1\mid T^{(j+1)}(n),
\]
it follows that
\[
T^{(j+1)}(n)\ge q_1\ge2.
\]
This completes the induction.

In particular,
\[
q_1\mid T^{(j)}(n),
\qquad
0\le j\le K.
\]
Since
\[
q_1\equiv-1\pmod M,
\]
we have
\[
M\mid q_1+1.
\]
For $0\le j<K$, write
\[
q_1^a\parallel T^{(j)}(n).
\]
Then
\[
\psi(q_1^a)
=
q_1^{a-1}(q_1+1)
\]
is divisible by $M$. By the multiplicativity of $\psi$, it follows
that
\[
M\mid\psi\bigl(T^{(j)}(n)\bigr).
\]
Therefore
\[
\begin{aligned}
T^{(j+1)}(n)
&=
\psi\bigl(T^{(j)}(n)\bigr)-T^{(j)}(n)\\
&\equiv
-T^{(j)}(n)
\pmod M.
\end{aligned}
\]
Starting from
\[
T^{(0)}(n)=n,
\]
induction gives
\[
T^{(j)}(n)\equiv(-1)^jn\pmod M,
\qquad
0\le j\le K.
\]
\end{proof}

\end{document}
